% !PRINTED 734
\weberpage{734}

%<<CONT
und hieraus ergiebt sich, dass $S'_\beta$ ebenso wie $S_\beta$ mit allen seinen conjugirten Werthen dasselbe Zeichen hat. $S'_\beta$ entsteht aber aus $S_\beta$ dadurch, dass $\mu$ an Stelle von $m$ gesetzt wird, und nach unserer Voraussetzung ist daher

\[
e_0 \equiv 0,\ e_1 \equiv 0,\ \ldots,\ e_{\nu''-1} \equiv 0 \ (\mathrm{mod}\ 2),
\]

und mit Hülfe von (19)

\[
e_{\nu''} \equiv 0,\ e_{\nu''+1} \equiv 0,\ \ldots,\ e_{\nu'-1} \equiv 0 \ (\mathrm{mod}\ 2),
\]

wodurch der Beweis des Satzes 2. geführt ist.

Wenn wir nun in der Formel (12) für $\mathcal{E}_0, \mathcal{E}_1, \ldots, \mathcal{E}_{\nu'-1}$ das System $\tau_0, \tau_1, \ldots, \tau_{\nu'-1}$ wählen und die Exponenten $\xi$ in die Form setzen:

\[
\xi_0 = \frac{e_0}{e},\ \xi_1 = \frac{e_1}{e},\ \cdots,\ \xi_{\nu'-1} = \frac{e_{\nu'-1}}{e},
\]

worin $e, e_0, e_1, \ldots, e_{\nu'-1}$ ganze rationale Zahlen ohne gemeinsamen Theiler sind, so erhalten wir

\[
\mathcal{E}(r)^e = \pm\, \tau_0^{e_0}\tau_1^{e_1} \cdots \tau_{\nu'-1}^{e_{\nu'-1}},
\]

und es folgt aus unserem Satze, dass $e$ \gesperrt{ungerade} sein muss; denn wäre es gerade, so wäre $\mathcal{E}(r)^e$ ein Quadrat einer reellen Grösse, und folglich hätte $\pm\, \mathcal{E}(r)^e$ mit seinen Conjugirten dasselbe Vorzeichen. Dann aber müssten nach unserem Satze alle $e_0, e_1, \ldots, e_{\nu'-1}$ durch 2 theilbar sein, was wieder gegen die Voraussetzung ist, dass die Exponenten $e, e_0, \ldots, e_{\nu'-1}$ keinen gemeinsamen Theiler haben sollen. Wir haben daher den Satz:

3. \gesperrt{Ist $\mathcal{E}(r)$ irgend eine Normaleinheit des Körpers $H_m$, so lassen sich die ungerade Zahl $e$ und die ganzen Zahlen $e_0, e_1, \ldots, e_{\nu'-1}$ so bestimmen, dass}

\weq{20}{\mathcal{E}(r_\beta)^e = \pm\, \tau_\beta^{e_0}\tau_{\beta+1}^{e_1} \cdots \tau_{\beta+\nu'-1}^{e_{\nu'-1}}}

\gesperrt{wird.}

Dies Ergebniss führt zu einigen wichtigen Consequenzen:

4. \gesperrt{Ist $T$ der Regulator des Systems der $\tau_\beta$, und $T_0$ der Regulator eines Fundamentalsystems normaler Einheiten in $H_m$, d.\,h. der Minimalwerth, den der Regulator irgend eines Systems unabhängiger Normaleinheiten annehmen kann, so ist}

\weq{21}{T = C\,T_0, \quad C \equiv 1\ (\mathrm{mod}\ 2),}

\gesperrt{worin $C$ eine ungerade natürliche Zahl bedeutet.}
